\scanpage{39}
\mspage{7}
où $I=D$ est le groupe des auto. de $a$, et où
les seules autres flèches, à part les $\gamma\in D$ et
l'identité de $b$, est une flèche $a\longrightarrow b$ (de
sorte que $b$ est objet final de $C$). Comme
$C$ admet un objet final, $\unsure{\widehat C|\widehat{\mathbb C}}$ \deleted{\unsure{est|et}} (donc aussi
$\unsure{\mathcal F|\mathfrak F}$) est de type d'homotopie trivial ($\infty$-\unsure{connexe|conn}),
ce qui signifie bien que $\unsure{\mathcal F|\mathfrak F}$ a le type d'
homotopie de $\deleted{\unsure{\mathbb C|C}}X=\mathbb C$.

Prenons maintenant le cas où $X$ est le
compactifié $\mathbb P^1_{\mathbb C}$ de $\mathbb C$ par un pt à l'infini,
\deleted{\unsure{donc}}\inserted{disons}
\[
\left\{\begin{array}{l}
X=\mathbb P^1_{\mathbb C}\simeq S^2\\
Y=\{\infty\}
\end{array}\right.
\]
\[
X^*=\deleted{\unsure{\mathbb P}}\mathbb C
\]
le diagramme (7) devient

\noindent (23)
\diagram{\begin{tikzcd}[row sep=large]
\mathbb D^*_{\infty} \arrow[r,hook,no head] \arrow[d] & \mathbb D_{\infty} \arrow[d,dotted] \\
\mathbb C \arrow[r,dashed] & \mathbb P^1_{\mathbb C}\simeq S^2
\end{tikzcd}}

dont la somme amalgamée est $S^2$, dont
le type d'homotopie n'a rien de trivial !
Le topos $\unsure{\mathcal F|\mathfrak F}$ se décrit maintenant à
l'aide du diagramme correspondant sur \unsure{les|le} $\Pi_1$,
savoir
% Pale foliation at the bottom left, in another hand.
\margin{38}
